\section{Cross-Resolution Supervision}
\label{sec:crossres}

\subsection{Student Architecture and Data}

Our student is a six-stage 3D residual-encoder nnU-Net with feature counts
$(32,64,128,256,320,320)$, one CT input channel, and two output classes.
It has $102{,}349{,}770$ trainable parameters.
The frozen training schedule contains 4,096 deterministic $192^3$ patches from
PHerc0139 and a single native-fine teacher record. We traverse this schedule
twice for 8,192 samples.
The small corpus reflects the time available before this progress submission,
not the extent of the registered data. Our broader pair plan contains three
additional non-prize training scrolls---PHerc0332, PHerc1667, and PHercParis4---
with 62,500 candidate origins allocated to each. None of the four training
scrolls is eligible for a Herculaneum prize.

Our held-out validation manifest contains registered PHerc0814 and PHerc1451
patches. We evaluate every 1,024 samples and retain model-only checkpoints at 1,024,
2,048, 4,096, and 8,192 samples. Because the measured duration response is not monotonic
(Sec.~\ref{sec:results}), we do not assume the final checkpoint is best.

\subsection{Fine-Teacher Occupancy and Medial Geometry}

The fine teacher provides two complementary targets in the student's physical
frame. The first is soft occupancy $q\in[0,1]$, obtained by an anti-aliased
pullback from the $2.399\,\mu$m grid. It estimates how much of a $9.362\,\mu$m
voxel is occupied by papyrus. The second is a binary medial crest $c$, projected
with an OR/max operator so that the existence of a fine-scale center survives
downsampling. Occupancy and crest have independent validity masks.

This separation follows a concrete failure of occupancy-only supervision. A
thin but real continuation commonly occupies only 25--45\% of a coarse voxel.
Soft cross-entropy is correctly calibrated when it drives such a voxel toward
$p=0.25$--$0.45$, but the result may remain below the segmentation threshold.
Reducing background separation moved both the center and its shoulders upward,
recovering length while reintroducing thick blobs. We therefore retain $q$ as
an amount-of-material target and use $c$ as an independent sheet-existence
target.

\paragraph{LSMAT-style center--radius representation.}
LSMAT represents a shape using medial centers and their maximally inscribed
radii rather than boundary samples alone~\cite{rebain2019lsmat}. We use that
center--radius view while retaining the released Villa teacher's center
construction. For each native-fine axial slice $T_z$, the teacher skeletonizes
the slice in 2D; the stacked centerlines receive one 3D morphological closing
and are intersected with the original foreground. We pair every surviving
center $u\in C_z$ with a physical 2D Euclidean-distance radius
\begin{equation}
 r_z(u)=\min_{y\notin T_z}\lVert u-y\rVert_{\mathrm{phys}},
 \qquad
 \widehat T_z=\bigcup_{u\in C_z} B^{\circ}(u,r_z(u)).
 \label{eq:center-radius}
\end{equation}
The disks are open because the stored distance reaches the nearest background
voxel center. Slice-wise radii are deliberate: pairing 2D centers with a 3D
distance transform would mix disks and spheres and collapse the desired medial
surface toward a curve. We audit reconstruction, component count, false-positive
girth, and halo stability before using these centers for supervision.

The representation gives us two independent geometric coordinates. Moving
along $C_z$ extends the papyrus sheet tangentially; increasing $r_z$ grows it
radially toward another wrap. That distinction is the reason for introducing
the LSMAT view: the connectivity loss should act along the first coordinate,
while the separation loss and radius screens police the second.

\paragraph{From medial centers to pins.}
We construct events only on training slices with at least 95\% valid teacher
support. Components smaller than 12 voxels are removed. A slice qualifies when
one thin teacher component $H_e$ contacts at least two disconnected M7
components $M_{e,i}$ and contains at least two teacher voxels absent from the
one-voxel M7 neighborhood. Both teacher and M7 must have maximum in-plane EDT
radius at most three voxels, and no more than 2\% of their foreground may
survive a two-voxel erosion; these checks reject blobs before they can define a
positive connectivity target.

Let $C_e$ be the projected medial centers inside $H_e$, and let $D_r$ denote
radius-$r$ dilation. We choose the smallest $d\in\{0,1,2,3\}$ for which
\begin{equation}
 \Omega_e = H_e\cap D_d(C_e),
 \qquad
 A_{e,i}=\Omega_e\cap D_2(M_{e,i})
 \label{eq:pins}
\end{equation}
contains every nonempty contact set $A_{e,i}$ in one connected component. The
sets $A_{e,i}$ are the \emph{pins}: they are where distinct released-M7
fragments meet the LSMAT-style teacher-medial corridor, not arbitrary endpoints
chosen from a student prediction. Corridor voxels within one voxel of M7, plus
all pin voxels, form the free-anchor set $F_e$. Events touching their owning
tile boundary are discarded, so every retained corridor is transformed and
loaded as one unit.

\subsection{Objective}

Let $p$ be the student's surface probability, $q$ the projected fine-teacher
occupancy, $c$ its binary medial-crest field,
$C=\{x:c_x=1\}$ the corresponding valid crest set, and $p_0$ the released M7
probability.
Our loss is
\begin{align}
\mathcal{L} ={}& \mathcal{L}_{\mathrm{CE}}(p,q)
 + 0.25\mathcal{L}_{\mathrm{Dice}}(p,q)
 + \mathcal{L}_{\mathrm{crest}}(p,c) \nonumber\\
&+2\mathcal{L}_{\mathrm{sep}}(p,q)
 +0.5\mathcal{L}_{\mathrm{KL}}(p,p_0)
 +\mathcal{L}_{\mathrm{pres}}(p,p_0) \nonumber\\
&+0.03125\mathcal{L}_{\mathrm{conn}}(p).
\label{eq:objective}
\end{align}
All terms are evaluated only where their corresponding validity fields permit.

\paragraph{Soft Occupancy and Crest Recall.}
Our cross-entropy and Dice targets retain fractional occupancy created by
fine-to-coarse projection. For a nonempty valid crest $C$, the additional term is
\begin{equation}
 \mathcal{L}_{\mathrm{crest}}
 = 1-\frac{\sum_{x\in C}p_x+1}{|C|+1}.
 \label{eq:crest-loss}
\end{equation}
We average Eq.~\ref{eq:crest-loss} over crest-bearing samples rather than over
all voxels. The normalization prevents sparse medial evidence from disappearing
into the background count, and the additive one is the smoothing constant used
by the released teacher loss. Because only crest voxels receive its gradient,
the term asks for longitudinal sheet evidence without defining a thick target.

\paragraph{De-Blob Separation.}
Thick predictions need explicit shrink pressure. An M7 blob can overlap the
fine sheet and still score well against an anti-aliased occupancy target.
Let $V$ and $V_c$ denote occupancy and crest validity. We form the thin positive
seed
\begin{equation}
 P=C\cup\bigl(\{q\geq0.5\}\cap(V\setminus V_c)\bigr),
 \label{eq:separation-positive}
\end{equation}
using the crest wherever it is known and falling back voxelwise to hard
occupancy only where it is not. We dilate $P$ by two voxels and retain only
occupancy-valid locations where the teacher is confidently background:
\begin{equation}
  S = \bigl(\operatorname{dilate}_2(P)\setminus P\bigr)
      \cap V\cap \{q\leq0.1\}.
  \label{eq:separation-shell}
\end{equation}
The separation loss is the mean background cross-entropy
\begin{equation}
  \mathcal{L}_{\mathrm{sep}}
  = -\frac{1}{|S|}\sum_{x\in S}\log(1-p_x).
  \label{eq:separation-loss}
\end{equation}
We weight this term by 2 in Eq.~\ref{eq:objective}. Because $S$ is a local fence around supported medial
structure rather than arbitrary background, the term attacks thickness and
short cross-wrap welds while the soft occupancy and crest terms remain free to
grow a faint sheet elsewhere.

\paragraph{M7 function and positive preservation.}
The teacher is locally more informative, not globally complete. We therefore
freeze a copy of released M7 during training and use its probability $p_0$ in
two differently masked terms. Let
$H=C\cup(\{q\geq0.5\}\cap V)$, $h_x=\mathbf{1}[x\in H]$, and
$m_x=\mathbf{1}[p_{0,x}\geq0.5]$. The KL mask is
\begin{align}
 K={}&\bigl((\overline V)\setminus D_2(H)\bigr) \nonumber\\
 &{}\cup\{x\in V:h_x=m_x\ \wedge\ (q_x\leq0.1\ \vee\ h_x=1)\}.
 \label{eq:kl-mask}
\end{align}
Thus unknown space is anchored except for a radius-two corridor around known
teacher positives. In known space, M7 is anchored only where its hard decision
agrees with a confident teacher decision; ambiguous partial occupancy is not
misclassified as agreement on background. We use
\begin{equation}
 \mathcal{L}_{\mathrm{KL}}
 =\frac{1}{|K|}\sum_{x\in K}
 D_{\mathrm{KL}}\!\left(\operatorname{Ber}(p_{0,x})\,
 \middle\|\,\operatorname{Ber}(p_x)\right).
 \label{eq:m7-kl}
\end{equation}

KL preservation alone can still allow a correct M7 foreground fragment to
shrink under the separation term. We therefore define
$R=V\cap D_2(H)\cap\{p_0\geq0.5\}$ and add the positive-only loss
\begin{equation}
 \mathcal{L}_{\mathrm{pres}}
 =-\frac{1}{|R|}\sum_{x\in R}\log p_x.
 \label{eq:m7-preservation}
\end{equation}
The two terms have separate roles: Eq.~\ref{eq:m7-kl} retains M7's foreground
and background function away from justified edits, while
Eq.~\ref{eq:m7-preservation} prevents de-blobbing from deleting M7-supported
sheet near the teacher. Teacher disagreement remains free to correct M7.

\paragraph{Dynamic medial connectivity.}
Separation supplies the desired shrink pressure, but without a complementary
topological term it can also sever the same faint longitudinal connections we
want to retain. Our first LSMAT axial loss connected the same pins with a
minimum-off-axis spanning tree, then raised the weakest 10\% of voxels on that
fixed route toward probability 0.2. This proved that a thin, center-biased route
could be supervised without using held-out geometry, but it also froze one of
many equally valid paths. The lowest weight only tied the incumbent; larger
weights monotonically reduced ordinary validation quality, and the strongest
weight also reduced anti-blob passes. The final loss therefore retains the
audited LSMAT corridor and pins but makes the path existential.

For event $e$, define the path capacity
\begin{equation}
 a_e(x)=
 \begin{cases}
 1, & x\in F_e,\\
 p_x, & x\in\Omega_e\setminus F_e,\\
 0, & x\notin\Omega_e.
 \end{cases}
 \label{eq:path-capacity}
\end{equation}
Taking $A_{e,0}$ as a source pin, initialize
$r_e^0(x)=\mathbf{1}[x\in A_{e,0}]$ and iterate
\begin{align}
 r_e^{t+1}(x)=\mathbf{1}[x\in\Omega_e]\max\!\Bigl(&r_e^t(x),
 \min\bigl(a_e(x),\max_{y\in\mathcal{N}_{26}(x)}r_e^t(y)\bigr)\Bigr).
 \label{eq:widest-recurrence}
\end{align}
This max--min recurrence computes the greatest bottleneck capacity among paths
of at most $t$ steps. After 96 steps, the event score and loss are
\begin{align}
 b_e & = \min_{i>0}\max_{x\in A_{e,i}} r_e^{96}(x), \nonumber\\
 \mathcal{L}_{\mathrm{conn}}
 & = \frac{1}{|E|}\sum_{e\in E}
 \max\bigl(0,\operatorname{logit}(0.2)-\operatorname{logit}(b_e)\bigr).
 \label{eq:connectivity-loss}
\end{align}
Each event receives one vote regardless of corridor size or pin count. Max/min
back-propagation places pressure on the current best path's limiting capacity,
not on every corridor voxel. Free M7 anchors have unit capacity and no direct
gradient; no connectivity gradient exists outside $\Omega_e$. The loss becomes
exactly zero when every pin is reachable through some path with bottleneck at
least 0.2.

The frozen atlas contains 149 fully owned events. Ninety-nine rows in the exact
schedule encounter an event, 477 rows are eligible, and the longest audited
event needs 44 of the allocated 96 propagation steps. Of these events, 135 have
two pins, 13 have three, and one has four; 122 need no crest dilation and 27 need
one voxel. The loss is evaluated only at the full-resolution output and has
weight 0.03125. Together, the separation shell and existential path constraint
express the two sides of de-blobbing: remove unsupported girth without deleting
sheet existence.

\paragraph{Deep supervision.}
The occupancy, crest, separation, KL, and preservation terms are applied at the
standard nnU-Net decoder outputs with normalized weights proportional to
$1,1/2,1/4,\ldots$ and zero weight on the coarsest output. Soft occupancy is
resized trilinearly. The sparse crest is max-pooled, while its validity is true
only if every contributing fine cell is valid. Connectivity is intentionally
excluded from lower-resolution outputs so resampling cannot broaden its medial
corridor.

\begin{figure*}[t]
  \centering
  \figureasset{figure_crossres_supervision}{0.96\linewidth}{2.0in}
  \caption{Actual PHerc0139 training data at dynamic event 138. Projected soft
  occupancy, its sparse medial crest, and the radius-two separation shell encode
  center/radius geometry without rewarding girth. The lower panels show the M7
  fragments, teacher-medial corridor, component-contact pins, and free anchors
  supplied to the max--min connectivity recurrence.}
  \Description{Training slice with projected teacher occupancy, medial crest,
  de-blob shell, M7 fragments, and dynamic-connectivity metadata.}
  \label{fig:supervision}
\end{figure*}

\subsection{Optimization and Trust Region}

We use Adam with constant learning rate $2\times10^{-5}$, betas $(0.9,0.999)$,
$\epsilon=10^{-8}$, and zero weight decay.
The batch size is three with eight-step gradient accumulation and bfloat16
autocast.
Seed 1203 fixes the sample order and augmentation sequence.

After each optimizer update, we project the complete parameter vector $\theta$
into a relative-$L^2$ ball around released parameters $\theta_0$:
\begin{equation}
\theta \leftarrow \theta_0 +
\min\!\left(1,\frac{r\lVert\theta_0\rVert_2}
{\lVert\theta-\theta_0\rVert_2}\right)(\theta-\theta_0),
\qquad r=0.0027535422.
\label{eq:trust}
\end{equation}
The projection is active on essentially every observed update.
It is thus part of the effective method rather than an inactive guardrail.
